%% kapa: exam=1 round=28 year=2017 period=1 form=A subject=real_estate subject_ko=부동산학원론 seq=16 number=96
%% source: Q-Net 감정평가사 기출문제 (cst00309) artlSeq=5199603 file=extracted/1008/5199603/2017년 제28회 감정평가사 1차시험 문제지/1교시 A형.pdf
%% conversion: kordoc 4.24.51 -> tools/gen1.py (auto); review=proofread-sonnet; flags=none
\begin{question}{96}
토지의 자연적ㆍ인문적 특성에 관한 설명으로 옳지 않은 것은?
\begin{choices}
  \item 부동성(위치의 고정성)으로 인해 외부효과가 발생한다.
  \item 분할ㆍ합병의 가능성은 용도의 다양성을 지원하는 특성이 있다.
  \item 용도의 다양성은 토지용도 중에서 최유효이용을 선택할 수 있는 근거가 된다.
  \item 일반적으로 부증성은 집약적 토지이용과 가격급등 현상을 일으키기도 한다.
  \item 토지의 인문적 특성 중에서 도시계획의 변경, 공업단지의 지정 등은 위치의 가변성 중 사회적 위치가 변화하는 예이다.
\end{choices}
\end{question}
\answer{5}
